% Answers to Chapter 1, from lectures/L01/appendix/c_solutions.md; the self-test from
% lectures/L01/README.md.

\facitavsnitt{c1}{Kapitel 1}{Talsystem och grundläggande aritmetik}
\begin{facit}

\svar{1.1} \sv{a} $13$ \sv{b} $25$ \sv{c} $5$ \sv{d} $15$

\svar{1.2} \sv{a} $\frac{5}{12}$ \sv{b} $\frac{7}{20}$ \sv{c} $\frac{3}{2} = 1{,}5$ \sv{d} $2$

\svar{1.3} \sv{a} $100$ \sv{b} $25\,\%$ \sv{c} Lägst $95\,\Omega$, högst $105\,\Omega$.

\svar{2.1} \sv{a} $\frac{1}{4}\,\Omega^{-1}$ \sv{b} $4\,\Omega$

\svar{2.2} \sv{a} $R_{\text{p}} = 6\,\text{k}\Omega$ med båda formlerna.
\sv{b} $10\,\text{k}\Omega$
\sv{c} Parallellkopplingen fungerar som en parentes: dess värde måste vara känt innan $R_1$ kan
adderas. Att addera alla tre motstånden ger fel svar, $36\,\text{k}\Omega$.

\svar{2.3} \sv{a} $2\,\text{mA}$ \sv{b} $8\,\text{V}$ \sv{c} $12\,\text{V}$
\sv{d} $I_2 = 0{,}5\,\text{mA}$, $I_3 = 1{,}5\,\text{mA}$; summan är $2\,\text{mA} = I$.

\svar{2.4} \sv{a} $8\,\text{V}$ \sv{b} $5{,}3\,\text{V}$

\svar{2.5} \sv{a} $75\,\%$ \sv{b} $21{,}6\,\text{W}$

\svar{2.6} \sv{a} $\mathbb{N}$ \sv{b} $\mathbb{Z}$ \sv{c} $\mathbb{Q}$ \sv{d} Irrationellt
\sv{e} $\mathbb{Q}$ \sv{f} Irrationellt

\svar{3.1} \sv{a} $-13$ \sv{b} $0$ \sv{c} $-6$ \sv{d} $16$ \sv{e} $1$

\svar{3.2} \sv{a} $-\abs{4} < -2{,}5 < 0 < \abs{-3} < \frac{7}{2}$ \sv{b} $4$
\sv{c} Båda är $5$: absolutbeloppet av en differens är avståndet mellan talen på tallinjen.
\sv{d} $\abs{U_1 - U_2} = 0{,}3\,\text{V}$

\svar{3.3} $\frac{1}{4} = 0{,}25 = 25\,\%$;\enspace $\frac{3}{5} = 0{,}6 = 60\,\%$;\enspace
$\frac{1}{8} = 0{,}125 = 12{,}5\,\%$;\enspace $\frac{3}{8} = 0{,}375 = 37{,}5\,\%$

\svar{3.4} \sv{a} $\frac{1}{3}$ \sv{b} $1$ \sv{c} $1$ \sv{d} $2$
\sv{e} Parallellresistansen av två motstånd på $4\,\Omega$, som är $2\,\Omega$.

\svar{3.5} \sv{a} $25\,\%$ \sv{b} $15\,\%$
\sv{c} Nej. $1{,}10 \times 0{,}90 = 0{,}99$, så värdet blir $1\,\%$ lägre än från början.
\sv{d} $75\,\text{W}$

\svar{3.6} \sv{a} $4{,}23\,\text{k}\Omega$ till $5{,}17\,\text{k}\Omega$ \sv{b} Ja.
\sv{c} Cirka $4{,}3\,\%$
\sv{d} Nej. Tillåtet är $209$--$231\,\Omega$, och $208\,\Omega$ avviker cirka $5{,}5\,\%$.

\svar{3.7} \sv{a} Sant: $n = \frac{n}{1}$.
\sv{b} Falskt, till exempel $\frac{1}{2}$.
\sv{c} Falskt: $\sqrt{9} = 3$.
\sv{d} Falskt, till exempel $\sqrt{2} + (-\sqrt{2}) = 0$.
\sv{e} Sant: $0{,}333\ldots = \frac{1}{3}$.
\sv{f} Falskt, till exempel $\pi$ och $\sqrt{2}$.

\svar{3.8} \sv{a} $50\,\Omega$ \sv{b} $25\,\Omega$
\sv{c} $n$ lika termer ger $\frac{1}{R_{\text{p}}} = \frac{n}{R}$.
\sv{d} $5$ motstånd

\svar{3.9} \sv{a} $20\,\Omega$ \sv{b} $280\,\Omega$ \sv{c} Ja, $12\,\Omega < 20\,\Omega$.
\sv{d} Nej. En parallellkoppling blir alltid mindre än det minsta ingående motståndet, här
$30\,\Omega$.

\svar{3.10} Uppifrån: $I = 3\,\text{mA}$;\enspace $U = 11\,\text{V}$;\enspace
$R = 3\,\text{k}\Omega$;\enspace $I = 5\,\text{mA}$

\svar{3.11} \sv{a} $190\,\text{W}$ \sv{b} $152\,\text{W}$ \sv{c} $76\,\%$ \sv{d} $48\,\text{W}$

\svar[Testa dig själv]{T} \sv{1} Irrationella tal, och därmed $\mathbb{R}$: $5$ är ingen
kvadrat av ett heltal, så $\sqrt{5}$ kan inte skrivas som ett bråk.
\sv{2} $3\,\Omega$

\end{facit}
